\documentclass[aspectratio=169,10pt]{beamer}
\makeatletter\def\input@path{{../../estilo/}{../estilo/}{estilo/}}\makeatother
\usepackage{devmobile}
\modulo[Flutter]{11}{Rede e APIs}{HTTP, REST, JSON, pacotes \cod{http} e Dio, interceptadores, erros, tempo real com WebSocket e testes}

\begin{document}
\capa

\begin{frame}{Objetivos do módulo}
  \begin{conceito}[Ao final deste módulo você será capaz de]
    \begin{itemize}
      \item Explicar requisição, resposta, métodos e códigos de status do HTTP.
      \item Consumir uma API REST com o pacote \cod{http} e com \textbf{Dio}.
      \item Converter JSON em objetos Dart à mão e com \cod{json\_serializable}.
      \item Usar interceptadores para autenticação, log e renovação de token.
      \item Tratar falhas de rede com \cod{Result}, tempo limite e retentativas.
      \item Receber dados em tempo real com WebSocket e testar o código de rede.
    \end{itemize}
  \end{conceito}
\end{frame}

\sumario

% =============================================================================
%  Fundamentos compartilhados --- Módulo 11: redes de computadores para apps
%  Incluído pelas trilhas Kotlin e Flutter com % =============================================================================
%  Fundamentos compartilhados --- Módulo 11: redes de computadores para apps
%  Incluído pelas trilhas Kotlin e Flutter com % =============================================================================
%  Fundamentos compartilhados --- Módulo 11: redes de computadores para apps
%  Incluído pelas trilhas Kotlin e Flutter com \input{../fundamentos/f11-redes}
% =============================================================================
\section{Fundamentos: redes de computadores}

\begin{frame}{Redes em camadas: o modelo TCP/IP}
  \begin{columns}[c]
    \begin{column}{0.5\textwidth}
      \centering
      \begin{tikzpicture}[every node/.style={font=\scriptsize}]
        \foreach \y/\nome/\prot/\estilo in {2.4/Aplicação/{HTTP, MQTT, DNS, WebSocket}/caixaescura,1.6/Transporte/{TCP, UDP, QUIC}/caixapetroleo,0.8/Internet/{IP (endereços e rotas)}/caixa,0/Enlace e física/{Wi-Fi, 4G/5G, Ethernet}/caixa}{
          \node[\estilo,minimum width=1.9cm,minimum height=6mm,font=\tiny\bfseries] at (0,\y) {\nome};
          \node[caixa,minimum width=3.2cm,minimum height=6mm,fill=white,font=\tiny,anchor=west] at (1.05,\y) {\prot};
        }
        \draw[setaambar] (-1.25,2.55) -- node[rotulo,left,text width=1.1cm,align=right]{cada camada envolve a de cima} (-1.25,-0.15);
      \end{tikzpicture}
      \begin{historia}
        Vint Cerf e Bob Kahn publicaram o TCP/IP em 1974; a ARPANET o adotou em 1983. O modelo OSI, de sete camadas, é a referência teórica.
      \end{historia}
    \end{column}
    \begin{column}{0.48\textwidth}
      \begin{definicao}[Protocolo]
        Conjunto de regras que define o formato das mensagens e a ordem da conversa entre dois computadores.
      \end{definicao}
      \begin{conceito}[Encapsulamento]
        Cada camada acrescenta seu cabeçalho e entrega à de baixo: a mensagem HTTP vai dentro de um segmento TCP, que vai dentro de um pacote IP, que vai dentro de um quadro de Wi-Fi ou 4G.
      \end{conceito}
      \begin{dica}
        O app só programa a camada de aplicação; o sistema operacional cuida do resto. Mas conhecer as de baixo explica lentidão e falhas.
      \end{dica}
    \end{column}
  \end{columns}
\end{frame}

\begin{frame}{Endereços, portas, NAT e DNS}
  \begin{columns}[T]
    \begin{column}{0.5\textwidth}
      \begin{definicao}[Endereço IP e porta]
        O IP identifica a máquina (IPv4 tem 32 bits; IPv6, 128 bits). A porta identifica o programa: 443 para HTTPS, 8883 para MQTT seguro.
      \end{definicao}
      \begin{definicao}[DNS]
        Sistema hierárquico que traduz nomes em IPs: o resolvedor consulta a raiz, depois o domínio de topo (\cod{.dev}) e por fim o servidor responsável. As respostas ficam em cache.
      \end{definicao}
    \end{column}
    \begin{column}{0.48\textwidth}
      \centering
      \begin{tikzpicture}[every node/.style={font=\tiny}]
        \node (c) at (0,0) {\celular[0.25]{}};
        \node[caixapetroleo,font=\tiny,minimum width=1.5cm,minimum height=10mm] (n) at (1.8,0) {NAT da\\ operadora};
        \node[caixaescura,font=\tiny,minimum width=1.4cm,minimum height=10mm] (s) at (5.0,0) {servidor};
        \draw[seta] (c.east |- 0,0.25) -- (n.west |- 0,0.25);
        \draw[seta] (n.east |- 0,0.25) -- node[rotulo,above]{sai: ok} (s.west |- 0,0.25);
        \draw[setaambar,dashed] (s.west |- 0,-0.25) -- node[rotulo,below]{entra: bloqueada} (n.east |- 0,-0.25);
\end{tikzpicture}
      \raggedright
      \begin{pordentro}[Por que o servidor não ``chama'' o celular]
        O celular fica atrás de um NAT da operadora, com endereço privado: ele abre conexões, mas não recebe conexões novas. Por isso existem as notificações \emph{push} e as conexões mantidas abertas pelo próprio app.
      \end{pordentro}
    \end{column}
  \end{columns}
\end{frame}

\begin{frame}{TCP e UDP}
  \begin{columns}[c]
    \begin{column}{0.4\textwidth}
      \centering
      \begin{tikzpicture}[x=1cm,y=0.55cm,every node/.style={font=\tiny}]
        \node[rotulo] at (0,0.6) {\textbf{app}};
        \node[rotulo] at (3,0.6) {\textbf{servidor}};
        \draw[cinzamedio] (0,0.2) -- (0,-5.2); \draw[cinzamedio] (3,0.2) -- (3,-5.2);
        \draw[seta] (0,-0.2) -- node[rotulo,above,sloped]{SYN} (3,-1);
        \draw[setaambar] (3,-1.3) -- node[rotulo,above,sloped]{SYN-ACK} (0,-2.1);
        \draw[seta] (0,-2.4) -- node[rotulo,above,sloped]{ACK + dados} (3,-3.2);
        \draw[setaambar] (3,-3.6) -- node[rotulo,above,sloped]{resposta} (0,-4.4);
        \node[rotulo,text width=3.4cm] at (1.5,-5.6) {aperto de mão em três vias: uma ida e volta antes dos dados};
      \end{tikzpicture}
    \end{column}
    \begin{column}{0.58\textwidth}
      \scriptsize
      \begin{tabularx}{\linewidth}{@{}l L L@{}}
        \toprule
         & \cabecalho{TCP} & \cabecalho{UDP} \\
        \midrule
        Conexão & estabelecida antes dos dados & nenhuma: cada pacote é independente \\
        Garantias & entrega, ordem e controle de congestionamento & nenhuma: pacotes podem se perder ou chegar fora de ordem \\
        Custo & mais idas e voltas & mínimo \\
        Usos & HTTP/1.1, HTTP/2, MQTT & DNS, voz, jogos, QUIC \\
        \bottomrule
      \end{tabularx}
      \normalsize
      \begin{definicao}[Latência e largura de banda]
        \textbf{Latência}: tempo de ida e volta de uma mensagem (RTT). \textbf{Largura de banda}: quantos dados passam por segundo. Em apps, as requisições são pequenas: a latência pesa mais.
      \end{definicao}
      \begin{pordentro}[Bloqueio na fila]
        No TCP, um pacote perdido segura todos os seguintes até ser reenviado, mesmo que pertençam a outras requisições.
      \end{pordentro}
    \end{column}
  \end{columns}
\end{frame}

\begin{frame}{TLS: como o HTTPS protege a conversa}
  \begin{columns}[c]
    \begin{column}{0.46\textwidth}
      \centering
      \begin{tikzpicture}[x=1cm,y=0.55cm,every node/.style={font=\tiny}]
        \node[rotulo] at (0,0.6) {\textbf{app}};
        \node[rotulo] at (3.2,0.6) {\textbf{servidor}};
        \draw[cinzamedio] (0,0.2) -- (0,-4.4); \draw[cinzamedio] (3.2,0.2) -- (3.2,-4.4);
        \draw[seta] (0,-0.2) -- node[rotulo,above,sloped]{olá + parte da chave} (3.2,-1);
        \draw[setaambar] (3.2,-1.3) -- node[rotulo,above,sloped,text width=2.6cm,align=center]{olá + parte da chave\\ + certificado assinado} (0,-2.3);
        \node[caixaambar,font=\tiny,text width=1.7cm] at (-0.2,-3.1) {confere o certificado};
        \draw[seta,line width=1.4pt] (0,-3.9) -- node[rotulo,above]{dados cifrados} (3.2,-3.9);
        \node[rotulo,text width=3.4cm] at (1.6,-4.9) {TLS 1.3: uma ida e volta para combinar as chaves};
      \end{tikzpicture}
    \end{column}
    \begin{column}{0.52\textwidth}
      \begin{definicao}[As três garantias do TLS]
        \textbf{Confidencialidade}: ninguém no caminho lê. \textbf{Integridade}: ninguém altera sem ser notado. \textbf{Autenticidade}: o app fala mesmo com o servidor certo.
      \end{definicao}
      \begin{pordentro}[Duas criptografias juntas]
        A criptografia assimétrica (pares de chaves) é usada só para combinar uma chave secreta; os dados seguem com criptografia simétrica, muito mais rápida (AES ou ChaCha20).
      \end{pordentro}
      \begin{conceito}[Cadeia de confiança]
        O certificado do servidor é assinado por uma autoridade certificadora em que o sistema operacional já confia. Certificado vencido ou de outro domínio: conexão recusada.
      \end{conceito}
    \end{column}
  \end{columns}
\end{frame}

\begin{frame}{A evolução do HTTP}
  \scriptsize
  \begin{tabularx}{\textwidth}{@{}l l L L@{}}
    \toprule
    \cabecalho{Versão} & \cabecalho{Ano} & \cabecalho{Novidade} & \cabecalho{Efeito no app} \\
    \midrule
    HTTP/1.0 & 1996 & uma conexão TCP por requisição & lento: aperto de mão a cada pedido \\
    HTTP/1.1 & 1997 & conexão reaproveitada (\emph{keep-alive}) & menos apertos de mão, mas um pedido por vez em cada conexão \\
    HTTP/2 & 2015 & binário, vários pedidos simultâneos na mesma conexão, cabeçalhos comprimidos & muitas requisições pequenas ficam baratas \\
    HTTP/3 & 2022 & roda sobre QUIC (UDP), sem bloqueio na fila do TCP & a conexão sobrevive à troca do Wi-Fi para o 4G \\
    \bottomrule
  \end{tabularx}
  \normalsize
  \vspace{4pt}
  \begin{columns}[T]
    \begin{column}{0.49\textwidth}
      \begin{historia}
        Tim Berners-Lee criou o HTTP e a Web no CERN, entre 1989 e 1991. Três décadas depois, o mesmo modelo de requisição e resposta move quase todos os apps.
      \end{historia}
    \end{column}
    \begin{column}{0.49\textwidth}
      \begin{pordentro}[Quem fala o quê]
        OkHttp usa HTTP/2; o Cronet, motor de rede do Chrome, também HTTP/3. No Flutter, \cod{cronet\_http} e \cod{cupertino\_http} usam as pilhas nativas do Android e do iOS.
      \end{pordentro}
    \end{column}
  \end{columns}
\end{frame}

\begin{frame}{Cache HTTP e requisições condicionais}
  \begin{columns}[c]
    \begin{column}{0.48\textwidth}
      \centering
      \begin{tikzpicture}[x=1cm,y=0.55cm,every node/.style={font=\tiny}]
        \node[rotulo] at (0,0.6) {\textbf{app}};
        \node[rotulo] at (3.6,0.6) {\textbf{servidor}};
        \draw[cinzamedio] (0,0.2) -- (0,-4.6); \draw[cinzamedio] (3.6,0.2) -- (3.6,-4.6);
        \draw[seta] (0,-0.2) -- node[rotulo,above,sloped]{\cod{GET /bancadas}} (3.6,-0.8);
        \draw[setaambar] (3.6,-1.1) -- node[rotulo,above,sloped]{200 + \cod{ETag: "v7"}} (0,-1.7);
        \node[rotulo] at (1.8,-2.2) {\ldots minutos depois\ldots};
        \draw[seta] (0,-2.7) -- node[rotulo,above,sloped]{\cod{If-None-Match: "v7"}} (3.6,-3.3);
        \draw[setaambar] (3.6,-3.6) -- node[rotulo,above,sloped]{304 Not Modified (sem corpo)} (0,-4.2);
      \end{tikzpicture}
    \end{column}
    \begin{column}{0.5\textwidth}
      \begin{definicao}[\emph{ETag}]
        Identificador da versão de um recurso. Na próxima vez, o app pergunta ``mudou desde a v7?''; se não mudou, a resposta 304 vem vazia e o app usa a cópia que já tem.
      \end{definicao}
      \scriptsize
      \begin{tabularx}{\linewidth}{@{}l L@{}}
        \toprule
        \cabecalho{\cod{Cache-Control}} & \cabecalho{Significado} \\
        \midrule
        \cod{max-age=60} & pode reutilizar por 60\,s sem perguntar \\
        \cod{no-cache} & pode guardar, mas pergunte antes de usar \\
        \cod{no-store} & não guarde (dados sensíveis) \\
        \bottomrule
      \end{tabularx}
      \normalsize
      \begin{dica}
        OkHttp e Dio têm cache HTTP: configurá-lo economiza dados móveis e bateria.
      \end{dica}
    \end{column}
  \end{columns}
\end{frame}

\begin{frame}{Serialização: de objetos a bytes}
  \begin{columns}[T]
    \begin{column}{0.52\textwidth}
      \begin{definicao}[Serialização]
        Converter objetos em bytes para enviar ou guardar; desserializar é o caminho de volta. Respostas comprimidas (gzip, brotli) reduzem muito o JSON.
      \end{definicao}
      \scriptsize
      \begin{tabularx}{\linewidth}{@{}l L L@{}}
        \toprule
        \cabecalho{Formato} & \cabecalho{Características} & \cabecalho{Uso típico} \\
        \midrule
        JSON & texto legível, sem esquema obrigatório & APIs REST \\
        Protocol Buffers & binário, com esquema, menor e mais rápido & gRPC, dados volumosos \\
        CBOR, MessagePack & JSON em formato binário & IoT \\
        \bottomrule
      \end{tabularx}
      \normalsize
    \end{column}
    \begin{column}{0.46\textwidth}
      \begin{conceito}[Evolução sem quebrar]
        Apps antigos continuam instalados por meses. O servidor só acrescenta campos; o app ignora campos desconhecidos e tem valores padrão para os que faltarem. Mudanças incompatíveis vão para uma nova versão da API (\cod{/v2}).
      \end{conceito}
      \begin{atencao}[JSON grande trava a tela]
        Decodificar megabytes de JSON é trabalho de CPU: faça fora da \emph{thread} da interface (Módulo 06).
      \end{atencao}
    \end{column}
  \end{columns}
\end{frame}

\begin{frame}{Redes móveis: latência, energia e falhas}
  \begin{columns}[T]
    \begin{column}{0.5\textwidth}
      \begin{pordentro}[O rádio gasta mesmo parado]
        Depois de cada transferência, o rádio celular fica alguns segundos em modo de alta potência antes de voltar ao repouso. Várias requisições espalhadas gastam mais bateria que as mesmas requisições agrupadas.
      \end{pordentro}
      \scriptsize
      \begin{tabularx}{\linewidth}{@{}l L@{}}
        \toprule
        \cabecalho{Método} & \cabecalho{Repetir é seguro?} \\
        \midrule
        \cod{GET}, \cod{PUT}, \cod{DELETE} & sim: são idempotentes (repetir dá o mesmo resultado) \\
        \cod{POST} & não: repetir pode criar duas reservas \\
        \bottomrule
      \end{tabularx}
    \end{column}
    \begin{column}{0.48\textwidth}
      \begin{conceito}[Defesas contra falhas]
        \textbf{Tempo limite} em toda requisição; \textbf{novas tentativas} com espera exponencial e um pouco de aleatoriedade (\emph{jitter}); \textbf{disjuntor} que para de insistir num servidor fora do ar.
      \end{conceito}
      \begin{atencao}[``Conectado'' não é ``com internet'']
        Wi-Fi de hotel com página de login, rede fraca ou servidor fora do ar: a única prova de rede é a requisição que funcionou.
      \end{atencao}
    \end{column}
  \end{columns}
\end{frame}

\begin{frame}{Tempo real: \emph{polling}, SSE, WebSocket e \emph{push}}
  \scriptsize
  \begin{tabularx}{\textwidth}{@{}l L L L@{}}
    \toprule
    \cabecalho{Técnica} & \cabecalho{Como funciona} & \cabecalho{Custo} & \cabecalho{Quando usar} \\
    \midrule
    \emph{Polling} & o app pergunta de tempos em tempos & alto: muitas requisições vazias & dados que mudam pouco \\
    \emph{Long polling} & o servidor segura a resposta até haver novidade & médio & quando nada melhor está disponível \\
    SSE & o servidor envia eventos por uma resposta HTTP que não termina & baixo; só servidor $\rightarrow$ app & painéis e notícias ao vivo \\
    WebSocket & canal aberto nos dois sentidos & baixo por mensagem; conexão aberta & chat, leituras ao vivo \\
    MQTT & publicação e assinatura por tópicos & muito baixo & sensores e IoT (Módulo 13) \\
    \emph{Push} (FCM, APNs) & o sistema entrega mesmo com o app fechado & quase zero para o app & avisos e alertas (Módulo 12) \\
    \bottomrule
  \end{tabularx}
  \normalsize
  \vspace{4pt}
  \begin{columns}[T]
    \begin{column}{0.49\textwidth}
      \begin{conceito}[Uma conexão só, do sistema]
        O \emph{push} usa uma única conexão mantida pelo sistema para todos os apps: é o jeito mais econômico de receber avisos.
      \end{conceito}
    \end{column}
    \begin{column}{0.49\textwidth}
      \begin{atencao}
        Conexões abertas não sobrevivem ao segundo plano: o sistema as encerra. Para avisar um app fechado, só o \emph{push}.
      \end{atencao}
    \end{column}
  \end{columns}
\end{frame}

% =============================================================================
\section{Fundamentos: redes de computadores}

\begin{frame}{Redes em camadas: o modelo TCP/IP}
  \begin{columns}[c]
    \begin{column}{0.5\textwidth}
      \centering
      \begin{tikzpicture}[every node/.style={font=\scriptsize}]
        \foreach \y/\nome/\prot/\estilo in {2.4/Aplicação/{HTTP, MQTT, DNS, WebSocket}/caixaescura,1.6/Transporte/{TCP, UDP, QUIC}/caixapetroleo,0.8/Internet/{IP (endereços e rotas)}/caixa,0/Enlace e física/{Wi-Fi, 4G/5G, Ethernet}/caixa}{
          \node[\estilo,minimum width=1.9cm,minimum height=6mm,font=\tiny\bfseries] at (0,\y) {\nome};
          \node[caixa,minimum width=3.2cm,minimum height=6mm,fill=white,font=\tiny,anchor=west] at (1.05,\y) {\prot};
        }
        \draw[setaambar] (-1.25,2.55) -- node[rotulo,left,text width=1.1cm,align=right]{cada camada envolve a de cima} (-1.25,-0.15);
      \end{tikzpicture}
      \begin{historia}
        Vint Cerf e Bob Kahn publicaram o TCP/IP em 1974; a ARPANET o adotou em 1983. O modelo OSI, de sete camadas, é a referência teórica.
      \end{historia}
    \end{column}
    \begin{column}{0.48\textwidth}
      \begin{definicao}[Protocolo]
        Conjunto de regras que define o formato das mensagens e a ordem da conversa entre dois computadores.
      \end{definicao}
      \begin{conceito}[Encapsulamento]
        Cada camada acrescenta seu cabeçalho e entrega à de baixo: a mensagem HTTP vai dentro de um segmento TCP, que vai dentro de um pacote IP, que vai dentro de um quadro de Wi-Fi ou 4G.
      \end{conceito}
      \begin{dica}
        O app só programa a camada de aplicação; o sistema operacional cuida do resto. Mas conhecer as de baixo explica lentidão e falhas.
      \end{dica}
    \end{column}
  \end{columns}
\end{frame}

\begin{frame}{Endereços, portas, NAT e DNS}
  \begin{columns}[T]
    \begin{column}{0.5\textwidth}
      \begin{definicao}[Endereço IP e porta]
        O IP identifica a máquina (IPv4 tem 32 bits; IPv6, 128 bits). A porta identifica o programa: 443 para HTTPS, 8883 para MQTT seguro.
      \end{definicao}
      \begin{definicao}[DNS]
        Sistema hierárquico que traduz nomes em IPs: o resolvedor consulta a raiz, depois o domínio de topo (\cod{.dev}) e por fim o servidor responsável. As respostas ficam em cache.
      \end{definicao}
    \end{column}
    \begin{column}{0.48\textwidth}
      \centering
      \begin{tikzpicture}[every node/.style={font=\tiny}]
        \node (c) at (0,0) {\celular[0.25]{}};
        \node[caixapetroleo,font=\tiny,minimum width=1.5cm,minimum height=10mm] (n) at (1.8,0) {NAT da\\ operadora};
        \node[caixaescura,font=\tiny,minimum width=1.4cm,minimum height=10mm] (s) at (5.0,0) {servidor};
        \draw[seta] (c.east |- 0,0.25) -- (n.west |- 0,0.25);
        \draw[seta] (n.east |- 0,0.25) -- node[rotulo,above]{sai: ok} (s.west |- 0,0.25);
        \draw[setaambar,dashed] (s.west |- 0,-0.25) -- node[rotulo,below]{entra: bloqueada} (n.east |- 0,-0.25);
\end{tikzpicture}
      \raggedright
      \begin{pordentro}[Por que o servidor não ``chama'' o celular]
        O celular fica atrás de um NAT da operadora, com endereço privado: ele abre conexões, mas não recebe conexões novas. Por isso existem as notificações \emph{push} e as conexões mantidas abertas pelo próprio app.
      \end{pordentro}
    \end{column}
  \end{columns}
\end{frame}

\begin{frame}{TCP e UDP}
  \begin{columns}[c]
    \begin{column}{0.4\textwidth}
      \centering
      \begin{tikzpicture}[x=1cm,y=0.55cm,every node/.style={font=\tiny}]
        \node[rotulo] at (0,0.6) {\textbf{app}};
        \node[rotulo] at (3,0.6) {\textbf{servidor}};
        \draw[cinzamedio] (0,0.2) -- (0,-5.2); \draw[cinzamedio] (3,0.2) -- (3,-5.2);
        \draw[seta] (0,-0.2) -- node[rotulo,above,sloped]{SYN} (3,-1);
        \draw[setaambar] (3,-1.3) -- node[rotulo,above,sloped]{SYN-ACK} (0,-2.1);
        \draw[seta] (0,-2.4) -- node[rotulo,above,sloped]{ACK + dados} (3,-3.2);
        \draw[setaambar] (3,-3.6) -- node[rotulo,above,sloped]{resposta} (0,-4.4);
        \node[rotulo,text width=3.4cm] at (1.5,-5.6) {aperto de mão em três vias: uma ida e volta antes dos dados};
      \end{tikzpicture}
    \end{column}
    \begin{column}{0.58\textwidth}
      \scriptsize
      \begin{tabularx}{\linewidth}{@{}l L L@{}}
        \toprule
         & \cabecalho{TCP} & \cabecalho{UDP} \\
        \midrule
        Conexão & estabelecida antes dos dados & nenhuma: cada pacote é independente \\
        Garantias & entrega, ordem e controle de congestionamento & nenhuma: pacotes podem se perder ou chegar fora de ordem \\
        Custo & mais idas e voltas & mínimo \\
        Usos & HTTP/1.1, HTTP/2, MQTT & DNS, voz, jogos, QUIC \\
        \bottomrule
      \end{tabularx}
      \normalsize
      \begin{definicao}[Latência e largura de banda]
        \textbf{Latência}: tempo de ida e volta de uma mensagem (RTT). \textbf{Largura de banda}: quantos dados passam por segundo. Em apps, as requisições são pequenas: a latência pesa mais.
      \end{definicao}
      \begin{pordentro}[Bloqueio na fila]
        No TCP, um pacote perdido segura todos os seguintes até ser reenviado, mesmo que pertençam a outras requisições.
      \end{pordentro}
    \end{column}
  \end{columns}
\end{frame}

\begin{frame}{TLS: como o HTTPS protege a conversa}
  \begin{columns}[c]
    \begin{column}{0.46\textwidth}
      \centering
      \begin{tikzpicture}[x=1cm,y=0.55cm,every node/.style={font=\tiny}]
        \node[rotulo] at (0,0.6) {\textbf{app}};
        \node[rotulo] at (3.2,0.6) {\textbf{servidor}};
        \draw[cinzamedio] (0,0.2) -- (0,-4.4); \draw[cinzamedio] (3.2,0.2) -- (3.2,-4.4);
        \draw[seta] (0,-0.2) -- node[rotulo,above,sloped]{olá + parte da chave} (3.2,-1);
        \draw[setaambar] (3.2,-1.3) -- node[rotulo,above,sloped,text width=2.6cm,align=center]{olá + parte da chave\\ + certificado assinado} (0,-2.3);
        \node[caixaambar,font=\tiny,text width=1.7cm] at (-0.2,-3.1) {confere o certificado};
        \draw[seta,line width=1.4pt] (0,-3.9) -- node[rotulo,above]{dados cifrados} (3.2,-3.9);
        \node[rotulo,text width=3.4cm] at (1.6,-4.9) {TLS 1.3: uma ida e volta para combinar as chaves};
      \end{tikzpicture}
    \end{column}
    \begin{column}{0.52\textwidth}
      \begin{definicao}[As três garantias do TLS]
        \textbf{Confidencialidade}: ninguém no caminho lê. \textbf{Integridade}: ninguém altera sem ser notado. \textbf{Autenticidade}: o app fala mesmo com o servidor certo.
      \end{definicao}
      \begin{pordentro}[Duas criptografias juntas]
        A criptografia assimétrica (pares de chaves) é usada só para combinar uma chave secreta; os dados seguem com criptografia simétrica, muito mais rápida (AES ou ChaCha20).
      \end{pordentro}
      \begin{conceito}[Cadeia de confiança]
        O certificado do servidor é assinado por uma autoridade certificadora em que o sistema operacional já confia. Certificado vencido ou de outro domínio: conexão recusada.
      \end{conceito}
    \end{column}
  \end{columns}
\end{frame}

\begin{frame}{A evolução do HTTP}
  \scriptsize
  \begin{tabularx}{\textwidth}{@{}l l L L@{}}
    \toprule
    \cabecalho{Versão} & \cabecalho{Ano} & \cabecalho{Novidade} & \cabecalho{Efeito no app} \\
    \midrule
    HTTP/1.0 & 1996 & uma conexão TCP por requisição & lento: aperto de mão a cada pedido \\
    HTTP/1.1 & 1997 & conexão reaproveitada (\emph{keep-alive}) & menos apertos de mão, mas um pedido por vez em cada conexão \\
    HTTP/2 & 2015 & binário, vários pedidos simultâneos na mesma conexão, cabeçalhos comprimidos & muitas requisições pequenas ficam baratas \\
    HTTP/3 & 2022 & roda sobre QUIC (UDP), sem bloqueio na fila do TCP & a conexão sobrevive à troca do Wi-Fi para o 4G \\
    \bottomrule
  \end{tabularx}
  \normalsize
  \vspace{4pt}
  \begin{columns}[T]
    \begin{column}{0.49\textwidth}
      \begin{historia}
        Tim Berners-Lee criou o HTTP e a Web no CERN, entre 1989 e 1991. Três décadas depois, o mesmo modelo de requisição e resposta move quase todos os apps.
      \end{historia}
    \end{column}
    \begin{column}{0.49\textwidth}
      \begin{pordentro}[Quem fala o quê]
        OkHttp usa HTTP/2; o Cronet, motor de rede do Chrome, também HTTP/3. No Flutter, \cod{cronet\_http} e \cod{cupertino\_http} usam as pilhas nativas do Android e do iOS.
      \end{pordentro}
    \end{column}
  \end{columns}
\end{frame}

\begin{frame}{Cache HTTP e requisições condicionais}
  \begin{columns}[c]
    \begin{column}{0.48\textwidth}
      \centering
      \begin{tikzpicture}[x=1cm,y=0.55cm,every node/.style={font=\tiny}]
        \node[rotulo] at (0,0.6) {\textbf{app}};
        \node[rotulo] at (3.6,0.6) {\textbf{servidor}};
        \draw[cinzamedio] (0,0.2) -- (0,-4.6); \draw[cinzamedio] (3.6,0.2) -- (3.6,-4.6);
        \draw[seta] (0,-0.2) -- node[rotulo,above,sloped]{\cod{GET /bancadas}} (3.6,-0.8);
        \draw[setaambar] (3.6,-1.1) -- node[rotulo,above,sloped]{200 + \cod{ETag: "v7"}} (0,-1.7);
        \node[rotulo] at (1.8,-2.2) {\ldots minutos depois\ldots};
        \draw[seta] (0,-2.7) -- node[rotulo,above,sloped]{\cod{If-None-Match: "v7"}} (3.6,-3.3);
        \draw[setaambar] (3.6,-3.6) -- node[rotulo,above,sloped]{304 Not Modified (sem corpo)} (0,-4.2);
      \end{tikzpicture}
    \end{column}
    \begin{column}{0.5\textwidth}
      \begin{definicao}[\emph{ETag}]
        Identificador da versão de um recurso. Na próxima vez, o app pergunta ``mudou desde a v7?''; se não mudou, a resposta 304 vem vazia e o app usa a cópia que já tem.
      \end{definicao}
      \scriptsize
      \begin{tabularx}{\linewidth}{@{}l L@{}}
        \toprule
        \cabecalho{\cod{Cache-Control}} & \cabecalho{Significado} \\
        \midrule
        \cod{max-age=60} & pode reutilizar por 60\,s sem perguntar \\
        \cod{no-cache} & pode guardar, mas pergunte antes de usar \\
        \cod{no-store} & não guarde (dados sensíveis) \\
        \bottomrule
      \end{tabularx}
      \normalsize
      \begin{dica}
        OkHttp e Dio têm cache HTTP: configurá-lo economiza dados móveis e bateria.
      \end{dica}
    \end{column}
  \end{columns}
\end{frame}

\begin{frame}{Serialização: de objetos a bytes}
  \begin{columns}[T]
    \begin{column}{0.52\textwidth}
      \begin{definicao}[Serialização]
        Converter objetos em bytes para enviar ou guardar; desserializar é o caminho de volta. Respostas comprimidas (gzip, brotli) reduzem muito o JSON.
      \end{definicao}
      \scriptsize
      \begin{tabularx}{\linewidth}{@{}l L L@{}}
        \toprule
        \cabecalho{Formato} & \cabecalho{Características} & \cabecalho{Uso típico} \\
        \midrule
        JSON & texto legível, sem esquema obrigatório & APIs REST \\
        Protocol Buffers & binário, com esquema, menor e mais rápido & gRPC, dados volumosos \\
        CBOR, MessagePack & JSON em formato binário & IoT \\
        \bottomrule
      \end{tabularx}
      \normalsize
    \end{column}
    \begin{column}{0.46\textwidth}
      \begin{conceito}[Evolução sem quebrar]
        Apps antigos continuam instalados por meses. O servidor só acrescenta campos; o app ignora campos desconhecidos e tem valores padrão para os que faltarem. Mudanças incompatíveis vão para uma nova versão da API (\cod{/v2}).
      \end{conceito}
      \begin{atencao}[JSON grande trava a tela]
        Decodificar megabytes de JSON é trabalho de CPU: faça fora da \emph{thread} da interface (Módulo 06).
      \end{atencao}
    \end{column}
  \end{columns}
\end{frame}

\begin{frame}{Redes móveis: latência, energia e falhas}
  \begin{columns}[T]
    \begin{column}{0.5\textwidth}
      \begin{pordentro}[O rádio gasta mesmo parado]
        Depois de cada transferência, o rádio celular fica alguns segundos em modo de alta potência antes de voltar ao repouso. Várias requisições espalhadas gastam mais bateria que as mesmas requisições agrupadas.
      \end{pordentro}
      \scriptsize
      \begin{tabularx}{\linewidth}{@{}l L@{}}
        \toprule
        \cabecalho{Método} & \cabecalho{Repetir é seguro?} \\
        \midrule
        \cod{GET}, \cod{PUT}, \cod{DELETE} & sim: são idempotentes (repetir dá o mesmo resultado) \\
        \cod{POST} & não: repetir pode criar duas reservas \\
        \bottomrule
      \end{tabularx}
    \end{column}
    \begin{column}{0.48\textwidth}
      \begin{conceito}[Defesas contra falhas]
        \textbf{Tempo limite} em toda requisição; \textbf{novas tentativas} com espera exponencial e um pouco de aleatoriedade (\emph{jitter}); \textbf{disjuntor} que para de insistir num servidor fora do ar.
      \end{conceito}
      \begin{atencao}[``Conectado'' não é ``com internet'']
        Wi-Fi de hotel com página de login, rede fraca ou servidor fora do ar: a única prova de rede é a requisição que funcionou.
      \end{atencao}
    \end{column}
  \end{columns}
\end{frame}

\begin{frame}{Tempo real: \emph{polling}, SSE, WebSocket e \emph{push}}
  \scriptsize
  \begin{tabularx}{\textwidth}{@{}l L L L@{}}
    \toprule
    \cabecalho{Técnica} & \cabecalho{Como funciona} & \cabecalho{Custo} & \cabecalho{Quando usar} \\
    \midrule
    \emph{Polling} & o app pergunta de tempos em tempos & alto: muitas requisições vazias & dados que mudam pouco \\
    \emph{Long polling} & o servidor segura a resposta até haver novidade & médio & quando nada melhor está disponível \\
    SSE & o servidor envia eventos por uma resposta HTTP que não termina & baixo; só servidor $\rightarrow$ app & painéis e notícias ao vivo \\
    WebSocket & canal aberto nos dois sentidos & baixo por mensagem; conexão aberta & chat, leituras ao vivo \\
    MQTT & publicação e assinatura por tópicos & muito baixo & sensores e IoT (Módulo 13) \\
    \emph{Push} (FCM, APNs) & o sistema entrega mesmo com o app fechado & quase zero para o app & avisos e alertas (Módulo 12) \\
    \bottomrule
  \end{tabularx}
  \normalsize
  \vspace{4pt}
  \begin{columns}[T]
    \begin{column}{0.49\textwidth}
      \begin{conceito}[Uma conexão só, do sistema]
        O \emph{push} usa uma única conexão mantida pelo sistema para todos os apps: é o jeito mais econômico de receber avisos.
      \end{conceito}
    \end{column}
    \begin{column}{0.49\textwidth}
      \begin{atencao}
        Conexões abertas não sobrevivem ao segundo plano: o sistema as encerra. Para avisar um app fechado, só o \emph{push}.
      \end{atencao}
    \end{column}
  \end{columns}
\end{frame}

% =============================================================================
\section{Fundamentos: redes de computadores}

\begin{frame}{Redes em camadas: o modelo TCP/IP}
  \begin{columns}[c]
    \begin{column}{0.5\textwidth}
      \centering
      \begin{tikzpicture}[every node/.style={font=\scriptsize}]
        \foreach \y/\nome/\prot/\estilo in {2.4/Aplicação/{HTTP, MQTT, DNS, WebSocket}/caixaescura,1.6/Transporte/{TCP, UDP, QUIC}/caixapetroleo,0.8/Internet/{IP (endereços e rotas)}/caixa,0/Enlace e física/{Wi-Fi, 4G/5G, Ethernet}/caixa}{
          \node[\estilo,minimum width=1.9cm,minimum height=6mm,font=\tiny\bfseries] at (0,\y) {\nome};
          \node[caixa,minimum width=3.2cm,minimum height=6mm,fill=white,font=\tiny,anchor=west] at (1.05,\y) {\prot};
        }
        \draw[setaambar] (-1.25,2.55) -- node[rotulo,left,text width=1.1cm,align=right]{cada camada envolve a de cima} (-1.25,-0.15);
      \end{tikzpicture}
      \begin{historia}
        Vint Cerf e Bob Kahn publicaram o TCP/IP em 1974; a ARPANET o adotou em 1983. O modelo OSI, de sete camadas, é a referência teórica.
      \end{historia}
    \end{column}
    \begin{column}{0.48\textwidth}
      \begin{definicao}[Protocolo]
        Conjunto de regras que define o formato das mensagens e a ordem da conversa entre dois computadores.
      \end{definicao}
      \begin{conceito}[Encapsulamento]
        Cada camada acrescenta seu cabeçalho e entrega à de baixo: a mensagem HTTP vai dentro de um segmento TCP, que vai dentro de um pacote IP, que vai dentro de um quadro de Wi-Fi ou 4G.
      \end{conceito}
      \begin{dica}
        O app só programa a camada de aplicação; o sistema operacional cuida do resto. Mas conhecer as de baixo explica lentidão e falhas.
      \end{dica}
    \end{column}
  \end{columns}
\end{frame}

\begin{frame}{Endereços, portas, NAT e DNS}
  \begin{columns}[T]
    \begin{column}{0.5\textwidth}
      \begin{definicao}[Endereço IP e porta]
        O IP identifica a máquina (IPv4 tem 32 bits; IPv6, 128 bits). A porta identifica o programa: 443 para HTTPS, 8883 para MQTT seguro.
      \end{definicao}
      \begin{definicao}[DNS]
        Sistema hierárquico que traduz nomes em IPs: o resolvedor consulta a raiz, depois o domínio de topo (\cod{.dev}) e por fim o servidor responsável. As respostas ficam em cache.
      \end{definicao}
    \end{column}
    \begin{column}{0.48\textwidth}
      \centering
      \begin{tikzpicture}[every node/.style={font=\tiny}]
        \node (c) at (0,0) {\celular[0.25]{}};
        \node[caixapetroleo,font=\tiny,minimum width=1.5cm,minimum height=10mm] (n) at (1.8,0) {NAT da\\ operadora};
        \node[caixaescura,font=\tiny,minimum width=1.4cm,minimum height=10mm] (s) at (5.0,0) {servidor};
        \draw[seta] (c.east |- 0,0.25) -- (n.west |- 0,0.25);
        \draw[seta] (n.east |- 0,0.25) -- node[rotulo,above]{sai: ok} (s.west |- 0,0.25);
        \draw[setaambar,dashed] (s.west |- 0,-0.25) -- node[rotulo,below]{entra: bloqueada} (n.east |- 0,-0.25);
\end{tikzpicture}
      \raggedright
      \begin{pordentro}[Por que o servidor não ``chama'' o celular]
        O celular fica atrás de um NAT da operadora, com endereço privado: ele abre conexões, mas não recebe conexões novas. Por isso existem as notificações \emph{push} e as conexões mantidas abertas pelo próprio app.
      \end{pordentro}
    \end{column}
  \end{columns}
\end{frame}

\begin{frame}{TCP e UDP}
  \begin{columns}[c]
    \begin{column}{0.4\textwidth}
      \centering
      \begin{tikzpicture}[x=1cm,y=0.55cm,every node/.style={font=\tiny}]
        \node[rotulo] at (0,0.6) {\textbf{app}};
        \node[rotulo] at (3,0.6) {\textbf{servidor}};
        \draw[cinzamedio] (0,0.2) -- (0,-5.2); \draw[cinzamedio] (3,0.2) -- (3,-5.2);
        \draw[seta] (0,-0.2) -- node[rotulo,above,sloped]{SYN} (3,-1);
        \draw[setaambar] (3,-1.3) -- node[rotulo,above,sloped]{SYN-ACK} (0,-2.1);
        \draw[seta] (0,-2.4) -- node[rotulo,above,sloped]{ACK + dados} (3,-3.2);
        \draw[setaambar] (3,-3.6) -- node[rotulo,above,sloped]{resposta} (0,-4.4);
        \node[rotulo,text width=3.4cm] at (1.5,-5.6) {aperto de mão em três vias: uma ida e volta antes dos dados};
      \end{tikzpicture}
    \end{column}
    \begin{column}{0.58\textwidth}
      \scriptsize
      \begin{tabularx}{\linewidth}{@{}l L L@{}}
        \toprule
         & \cabecalho{TCP} & \cabecalho{UDP} \\
        \midrule
        Conexão & estabelecida antes dos dados & nenhuma: cada pacote é independente \\
        Garantias & entrega, ordem e controle de congestionamento & nenhuma: pacotes podem se perder ou chegar fora de ordem \\
        Custo & mais idas e voltas & mínimo \\
        Usos & HTTP/1.1, HTTP/2, MQTT & DNS, voz, jogos, QUIC \\
        \bottomrule
      \end{tabularx}
      \normalsize
      \begin{definicao}[Latência e largura de banda]
        \textbf{Latência}: tempo de ida e volta de uma mensagem (RTT). \textbf{Largura de banda}: quantos dados passam por segundo. Em apps, as requisições são pequenas: a latência pesa mais.
      \end{definicao}
      \begin{pordentro}[Bloqueio na fila]
        No TCP, um pacote perdido segura todos os seguintes até ser reenviado, mesmo que pertençam a outras requisições.
      \end{pordentro}
    \end{column}
  \end{columns}
\end{frame}

\begin{frame}{TLS: como o HTTPS protege a conversa}
  \begin{columns}[c]
    \begin{column}{0.46\textwidth}
      \centering
      \begin{tikzpicture}[x=1cm,y=0.55cm,every node/.style={font=\tiny}]
        \node[rotulo] at (0,0.6) {\textbf{app}};
        \node[rotulo] at (3.2,0.6) {\textbf{servidor}};
        \draw[cinzamedio] (0,0.2) -- (0,-4.4); \draw[cinzamedio] (3.2,0.2) -- (3.2,-4.4);
        \draw[seta] (0,-0.2) -- node[rotulo,above,sloped]{olá + parte da chave} (3.2,-1);
        \draw[setaambar] (3.2,-1.3) -- node[rotulo,above,sloped,text width=2.6cm,align=center]{olá + parte da chave\\ + certificado assinado} (0,-2.3);
        \node[caixaambar,font=\tiny,text width=1.7cm] at (-0.2,-3.1) {confere o certificado};
        \draw[seta,line width=1.4pt] (0,-3.9) -- node[rotulo,above]{dados cifrados} (3.2,-3.9);
        \node[rotulo,text width=3.4cm] at (1.6,-4.9) {TLS 1.3: uma ida e volta para combinar as chaves};
      \end{tikzpicture}
    \end{column}
    \begin{column}{0.52\textwidth}
      \begin{definicao}[As três garantias do TLS]
        \textbf{Confidencialidade}: ninguém no caminho lê. \textbf{Integridade}: ninguém altera sem ser notado. \textbf{Autenticidade}: o app fala mesmo com o servidor certo.
      \end{definicao}
      \begin{pordentro}[Duas criptografias juntas]
        A criptografia assimétrica (pares de chaves) é usada só para combinar uma chave secreta; os dados seguem com criptografia simétrica, muito mais rápida (AES ou ChaCha20).
      \end{pordentro}
      \begin{conceito}[Cadeia de confiança]
        O certificado do servidor é assinado por uma autoridade certificadora em que o sistema operacional já confia. Certificado vencido ou de outro domínio: conexão recusada.
      \end{conceito}
    \end{column}
  \end{columns}
\end{frame}

\begin{frame}{A evolução do HTTP}
  \scriptsize
  \begin{tabularx}{\textwidth}{@{}l l L L@{}}
    \toprule
    \cabecalho{Versão} & \cabecalho{Ano} & \cabecalho{Novidade} & \cabecalho{Efeito no app} \\
    \midrule
    HTTP/1.0 & 1996 & uma conexão TCP por requisição & lento: aperto de mão a cada pedido \\
    HTTP/1.1 & 1997 & conexão reaproveitada (\emph{keep-alive}) & menos apertos de mão, mas um pedido por vez em cada conexão \\
    HTTP/2 & 2015 & binário, vários pedidos simultâneos na mesma conexão, cabeçalhos comprimidos & muitas requisições pequenas ficam baratas \\
    HTTP/3 & 2022 & roda sobre QUIC (UDP), sem bloqueio na fila do TCP & a conexão sobrevive à troca do Wi-Fi para o 4G \\
    \bottomrule
  \end{tabularx}
  \normalsize
  \vspace{4pt}
  \begin{columns}[T]
    \begin{column}{0.49\textwidth}
      \begin{historia}
        Tim Berners-Lee criou o HTTP e a Web no CERN, entre 1989 e 1991. Três décadas depois, o mesmo modelo de requisição e resposta move quase todos os apps.
      \end{historia}
    \end{column}
    \begin{column}{0.49\textwidth}
      \begin{pordentro}[Quem fala o quê]
        OkHttp usa HTTP/2; o Cronet, motor de rede do Chrome, também HTTP/3. No Flutter, \cod{cronet\_http} e \cod{cupertino\_http} usam as pilhas nativas do Android e do iOS.
      \end{pordentro}
    \end{column}
  \end{columns}
\end{frame}

\begin{frame}{Cache HTTP e requisições condicionais}
  \begin{columns}[c]
    \begin{column}{0.48\textwidth}
      \centering
      \begin{tikzpicture}[x=1cm,y=0.55cm,every node/.style={font=\tiny}]
        \node[rotulo] at (0,0.6) {\textbf{app}};
        \node[rotulo] at (3.6,0.6) {\textbf{servidor}};
        \draw[cinzamedio] (0,0.2) -- (0,-4.6); \draw[cinzamedio] (3.6,0.2) -- (3.6,-4.6);
        \draw[seta] (0,-0.2) -- node[rotulo,above,sloped]{\cod{GET /bancadas}} (3.6,-0.8);
        \draw[setaambar] (3.6,-1.1) -- node[rotulo,above,sloped]{200 + \cod{ETag: "v7"}} (0,-1.7);
        \node[rotulo] at (1.8,-2.2) {\ldots minutos depois\ldots};
        \draw[seta] (0,-2.7) -- node[rotulo,above,sloped]{\cod{If-None-Match: "v7"}} (3.6,-3.3);
        \draw[setaambar] (3.6,-3.6) -- node[rotulo,above,sloped]{304 Not Modified (sem corpo)} (0,-4.2);
      \end{tikzpicture}
    \end{column}
    \begin{column}{0.5\textwidth}
      \begin{definicao}[\emph{ETag}]
        Identificador da versão de um recurso. Na próxima vez, o app pergunta ``mudou desde a v7?''; se não mudou, a resposta 304 vem vazia e o app usa a cópia que já tem.
      \end{definicao}
      \scriptsize
      \begin{tabularx}{\linewidth}{@{}l L@{}}
        \toprule
        \cabecalho{\cod{Cache-Control}} & \cabecalho{Significado} \\
        \midrule
        \cod{max-age=60} & pode reutilizar por 60\,s sem perguntar \\
        \cod{no-cache} & pode guardar, mas pergunte antes de usar \\
        \cod{no-store} & não guarde (dados sensíveis) \\
        \bottomrule
      \end{tabularx}
      \normalsize
      \begin{dica}
        OkHttp e Dio têm cache HTTP: configurá-lo economiza dados móveis e bateria.
      \end{dica}
    \end{column}
  \end{columns}
\end{frame}

\begin{frame}{Serialização: de objetos a bytes}
  \begin{columns}[T]
    \begin{column}{0.52\textwidth}
      \begin{definicao}[Serialização]
        Converter objetos em bytes para enviar ou guardar; desserializar é o caminho de volta. Respostas comprimidas (gzip, brotli) reduzem muito o JSON.
      \end{definicao}
      \scriptsize
      \begin{tabularx}{\linewidth}{@{}l L L@{}}
        \toprule
        \cabecalho{Formato} & \cabecalho{Características} & \cabecalho{Uso típico} \\
        \midrule
        JSON & texto legível, sem esquema obrigatório & APIs REST \\
        Protocol Buffers & binário, com esquema, menor e mais rápido & gRPC, dados volumosos \\
        CBOR, MessagePack & JSON em formato binário & IoT \\
        \bottomrule
      \end{tabularx}
      \normalsize
    \end{column}
    \begin{column}{0.46\textwidth}
      \begin{conceito}[Evolução sem quebrar]
        Apps antigos continuam instalados por meses. O servidor só acrescenta campos; o app ignora campos desconhecidos e tem valores padrão para os que faltarem. Mudanças incompatíveis vão para uma nova versão da API (\cod{/v2}).
      \end{conceito}
      \begin{atencao}[JSON grande trava a tela]
        Decodificar megabytes de JSON é trabalho de CPU: faça fora da \emph{thread} da interface (Módulo 06).
      \end{atencao}
    \end{column}
  \end{columns}
\end{frame}

\begin{frame}{Redes móveis: latência, energia e falhas}
  \begin{columns}[T]
    \begin{column}{0.5\textwidth}
      \begin{pordentro}[O rádio gasta mesmo parado]
        Depois de cada transferência, o rádio celular fica alguns segundos em modo de alta potência antes de voltar ao repouso. Várias requisições espalhadas gastam mais bateria que as mesmas requisições agrupadas.
      \end{pordentro}
      \scriptsize
      \begin{tabularx}{\linewidth}{@{}l L@{}}
        \toprule
        \cabecalho{Método} & \cabecalho{Repetir é seguro?} \\
        \midrule
        \cod{GET}, \cod{PUT}, \cod{DELETE} & sim: são idempotentes (repetir dá o mesmo resultado) \\
        \cod{POST} & não: repetir pode criar duas reservas \\
        \bottomrule
      \end{tabularx}
    \end{column}
    \begin{column}{0.48\textwidth}
      \begin{conceito}[Defesas contra falhas]
        \textbf{Tempo limite} em toda requisição; \textbf{novas tentativas} com espera exponencial e um pouco de aleatoriedade (\emph{jitter}); \textbf{disjuntor} que para de insistir num servidor fora do ar.
      \end{conceito}
      \begin{atencao}[``Conectado'' não é ``com internet'']
        Wi-Fi de hotel com página de login, rede fraca ou servidor fora do ar: a única prova de rede é a requisição que funcionou.
      \end{atencao}
    \end{column}
  \end{columns}
\end{frame}

\begin{frame}{Tempo real: \emph{polling}, SSE, WebSocket e \emph{push}}
  \scriptsize
  \begin{tabularx}{\textwidth}{@{}l L L L@{}}
    \toprule
    \cabecalho{Técnica} & \cabecalho{Como funciona} & \cabecalho{Custo} & \cabecalho{Quando usar} \\
    \midrule
    \emph{Polling} & o app pergunta de tempos em tempos & alto: muitas requisições vazias & dados que mudam pouco \\
    \emph{Long polling} & o servidor segura a resposta até haver novidade & médio & quando nada melhor está disponível \\
    SSE & o servidor envia eventos por uma resposta HTTP que não termina & baixo; só servidor $\rightarrow$ app & painéis e notícias ao vivo \\
    WebSocket & canal aberto nos dois sentidos & baixo por mensagem; conexão aberta & chat, leituras ao vivo \\
    MQTT & publicação e assinatura por tópicos & muito baixo & sensores e IoT (Módulo 13) \\
    \emph{Push} (FCM, APNs) & o sistema entrega mesmo com o app fechado & quase zero para o app & avisos e alertas (Módulo 12) \\
    \bottomrule
  \end{tabularx}
  \normalsize
  \vspace{4pt}
  \begin{columns}[T]
    \begin{column}{0.49\textwidth}
      \begin{conceito}[Uma conexão só, do sistema]
        O \emph{push} usa uma única conexão mantida pelo sistema para todos os apps: é o jeito mais econômico de receber avisos.
      \end{conceito}
    \end{column}
    \begin{column}{0.49\textwidth}
      \begin{atencao}
        Conexões abertas não sobrevivem ao segundo plano: o sistema as encerra. Para avisar um app fechado, só o \emph{push}.
      \end{atencao}
    \end{column}
  \end{columns}
\end{frame}


% =============================================================================
\section{Fundamentos de HTTP}

\begin{frame}{Requisição e resposta}
  \begin{columns}[c]
    \begin{column}{0.52\textwidth}
      \centering
      \begin{tikzpicture}
        \node[caixa,minimum width=1.8cm,minimum height=1.6cm] (app) {App\\ \tiny cliente};
        \node[caixaescura,minimum width=1.8cm,minimum height=1.6cm,right=3.2cm of app] (srv) {API\\ \tiny servidor};
        \draw[setaambar] ([yshift=4mm]app.east) -- node[rotulo,above,align=center]{\tiny GET /v1/bancadas\\ \tiny Authorization: Bearer \ldots} ([yshift=4mm]srv.west);
        \draw[seta] ([yshift=-4mm]srv.west) -- node[rotulo,below,align=center]{\tiny 200 OK\\ \tiny \texttt{[\{"id":1,"nome":"B1"\}]}} ([yshift=-4mm]app.east);
      \end{tikzpicture}
      \vspace{6pt}
      \begin{conceito}[Partes de uma requisição]
        \textbf{Método} + \textbf{URL} + \textbf{cabeçalhos} + \textbf{corpo} (opcional). A resposta traz \textbf{código de status}, cabeçalhos e corpo, em geral JSON.
      \end{conceito}
    \end{column}
    \begin{column}{0.46\textwidth}
      \scriptsize
      \begin{tabularx}{\linewidth}{@{}l L@{}}
        \toprule
        \cabecalho{Método} & \cabecalho{Uso em REST} \\
        \midrule
        \cod{GET} & ler um recurso ou uma lista \\
        \cod{POST} & criar um recurso \\
        \cod{PUT} & substituir um recurso inteiro \\
        \cod{PATCH} & alterar parte de um recurso \\
        \cod{DELETE} & remover um recurso \\
        \bottomrule
      \end{tabularx}
      \vspace{4pt}
      \begin{tabularx}{\linewidth}{@{}l L@{}}
        \toprule
        \cabecalho{Status} & \cabecalho{Significado} \\
        \midrule
        \cod{2xx} & sucesso (200 OK, 201 Created, 204) \\
        \cod{4xx} & erro do cliente (400, 401, 403, 404, 409, 422) \\
        \cod{5xx} & erro do servidor (500, 502, 503) \\
        \bottomrule
      \end{tabularx}
    \end{column}
  \end{columns}
\end{frame}

\begin{frame}{Estilos de API}
  \footnotesize
  \begin{tabularx}{\textwidth}{@{}l L L@{}}
    \toprule
    \cabecalho{Estilo} & \cabecalho{Como funciona} & \cabecalho{Quando usar} \\
    \midrule
    \textbf{REST} & recursos em URLs, métodos HTTP, JSON & padrão da maioria dos backends; foco deste módulo \\
    \textbf{GraphQL} & uma URL; o cliente descreve os campos que quer & telas que combinam muitos dados; evita várias chamadas \\
    \textbf{gRPC} & contratos Protocol Buffers sobre HTTP/2 & comunicação entre serviços; alto desempenho \\
    \textbf{WebSocket} & conexão aberta nos dois sentidos & tempo real: leituras ao vivo, chat \\
    \textbf{MQTT} & publicação e assinatura por tópicos & IoT com rede instável (Módulo 13) \\
    \bottomrule
  \end{tabularx}
  \vspace{4pt}
  \begin{columns}[T]
    \begin{column}{0.47\textwidth}
      \begin{dica}[Explore antes de codificar]
        Teste a API no \textbf{Bruno}, \textbf{Postman} ou \cod{curl}. Se ela publica um documento \textbf{OpenAPI}, leia os contratos ali.
      \end{dica}
    \end{column}
    \begin{column}{0.47\textwidth}
      \begin{atencao}
        Use sempre \textbf{HTTPS}. O Android bloqueia HTTP sem TLS por padrão e o iOS exige TLS pelo ATS.
      \end{atencao}
    \end{column}
  \end{columns}
\end{frame}

% =============================================================================
\section{O pacote http}

\begin{frame}[fragile]{Preparando o projeto}
  \begin{columns}[T]
    \begin{column}{0.5\textwidth}
\begin{terminal}
$ flutter pub add http
\end{terminal}
\begin{codigo}[title=android/app/src/main/AndroidManifest.xml]{XML}
<manifest ...>
  <uses-permission
      android:name="android.permission.INTERNET" />
  <application ...>
\end{codigo}
\begin{codigo}[title=macos/Runner/*.entitlements]{XML}
<key>com.apple.security.network.client</key>
<true/>
\end{codigo}
    \end{column}
    \begin{column}{0.48\textwidth}
      \begin{erro}[Funciona no debug, falha na loja]
        O modo debug do Android já inclui a permissão \cod{INTERNET}; a versão \emph{release} não. Sem a linha no manifesto, o app publicado não acessa a rede.
      \end{erro}
      \begin{conceito}[Plataformas]
        iOS e web não pedem permissão para rede. No macOS, o app roda em \emph{sandbox} e precisa do \emph{entitlement} de cliente de rede.
      \end{conceito}
    \end{column}
  \end{columns}
\end{frame}

\begin{frame}[fragile]{Um serviço de API com \cod{http}}
  \begin{columns}[T]
    \begin{column}{0.6\textwidth}
\begin{codigo}[title=data/services/api\_service.dart]{Dart}
class ApiService {
  ApiService({http.Client? client})
      : _c = client ?? http.Client();
  final http.Client _c;
  static const _base = 'https://api.labmonitor.dev/v1';

  Future<List<BancadaDto>> getBancadas() async {
    final r = await _c.get(Uri.parse('$_base/bancadas'))
        .timeout(const Duration(seconds: 10));
    if (r.statusCode != 200) throw ApiException(r.statusCode);
    final lista = jsonDecode(r.body) as List<dynamic>;
    return lista.map((j) => BancadaDto.fromJson(j)).toList();
  }
  Future<void> reservar(int id, DateTime ate) async {
    final r = await _c.post(
      Uri.parse('$_base/bancadas/$id/reservas'),
      headers: {'Content-Type': 'application/json'},
      body: jsonEncode({'ate': ate.toIso8601String()}));
    if (r.statusCode != 201) throw ApiException(r.statusCode);
  }
}
\end{codigo}
    \end{column}
    \begin{column}{0.38\textwidth}
      \begin{conceito}[Boas práticas]
        \begin{itemize}
          \item Um \cod{http.Client} reutilizado mantém as conexões abertas.
          \item Receber o cliente no construtor permite testar com \cod{MockClient}.
          \item Sempre defina um tempo limite.
          \item Confira o status antes de ler o corpo.
        \end{itemize}
      \end{conceito}
      \begin{dica}
        \cod{ApiException} é uma classe sua; não use \cod{HttpException} de \cod{dart:io}, que não existe na web.
      \end{dica}
    \end{column}
  \end{columns}
\end{frame}

% =============================================================================
\section{JSON}

\begin{frame}[fragile]{JSON à mão com \emph{patterns}}
  \begin{columns}[T]
    \begin{column}{0.56\textwidth}
\begin{codigo}[title=domain/dto/bancada\_dto.dart]{Dart}
class BancadaDto {
  BancadaDto(this.id, this.nome, this.reservadaAte);
  final int id;
  final String nome;
  final DateTime? reservadaAte;

  factory BancadaDto.fromJson(Object? json) {
    if (json case {
      'id': int id,
      'nome': String nome,
      'reservada_ate': String? ate,
    }) {
      return BancadaDto(id, nome,
          ate == null ? null : DateTime.parse(ate));
    }
    throw FormatException('Bancada inválida: $json');
  }

  Map<String, Object?> toJson() => {'id': id, 'nome': nome,
      'reservada_ate': reservadaAte?.toIso8601String()};
}
\end{codigo}
    \end{column}
    \begin{column}{0.42\textwidth}
      \begin{conceito}[\emph{Pattern} de mapa]
        O \cod{if case} confere a estrutura e os tipos e extrai os campos em uma só expressão (Dart 3, Módulo 02). JSON inesperado vira um erro claro, não um \cod{null} silencioso.
      \end{conceito}
      \begin{atencao}
        \cod{jsonDecode} devolve \cod{dynamic}. Acessar \cod{json['nome']} sem conferir o tipo adia o erro para longe da sua causa.
      \end{atencao}
    \end{column}
  \end{columns}
\end{frame}

\begin{frame}[fragile]{JSON com geração de código}
  \begin{columns}[T]
    \begin{column}{0.58\textwidth}
\begin{codigo}[title=Com json\_serializable]{Dart}
import 'package:json_annotation/json_annotation.dart';
part 'bancada_dto.g.dart';
@JsonSerializable(fieldRename: FieldRename.snake)
class BancadaDto {
  BancadaDto({required this.id, required this.nome,
      this.reservadaAte});
  final int id;
  final String nome;
  final DateTime? reservadaAte;    // "reservada_ate"
  factory BancadaDto.fromJson(Map<String, dynamic> json) =>
      _$BancadaDtoFromJson(json);
  Map<String, dynamic> toJson() => _$BancadaDtoToJson(this);
}
\end{codigo}
\begin{terminal}
$ flutter pub add json_annotation
$ flutter pub add -d json_serializable build_runner
$ dart run build_runner watch
\end{terminal}
    \end{column}
    \begin{column}{0.4\textwidth}
      \begin{conceito}[Quando gerar]
        Com poucos modelos, o código à mão basta. Com dezenas de modelos e campos aninhados, a geração evita erros de digitação.
      \end{conceito}
      \begin{dica}[\cod{freezed}]
        Gera também \cod{copyWith}, igualdade, \cod{toString} e uniões seladas; muito usado com Riverpod e Bloc.
      \end{dica}
      \begin{dica}[JSON grande]
        Decodifique fora da \emph{thread} da UI com \cod{Isolate.run} (Módulo 06).
      \end{dica}
    \end{column}
  \end{columns}
\end{frame}

% =============================================================================
\section{Dio}

\begin{frame}[fragile]{Configurando o Dio}
  \begin{columns}[T]
    \begin{column}{0.58\textwidth}
\begin{codigo}[title=data/services/cliente\_http.dart]{Dart}
Dio criarDio(TokenStorage tokens) {
  final dio = Dio(BaseOptions(
    baseUrl: 'https://api.labmonitor.dev/v1/',
    connectTimeout: const Duration(seconds: 5),
    receiveTimeout: const Duration(seconds: 10),
  ));
  dio.interceptors.addAll([
    InterceptorsWrapper(onRequest: (opcoes, handler) async {
      final t = await tokens.lerAcesso();
      if (t != null) opcoes.headers['Authorization'] = 'Bearer $t';
      handler.next(opcoes);
    }),
    if (kDebugMode) LogInterceptor(responseBody: true),
  ]);
  return dio;
}
// uso: o JSON já vem decodificado
final r = await dio.get<List<dynamic>>('bancadas',
    queryParameters: {'lab': 3, 'livres': true});
final bancadas =
    r.data!.map((j) => BancadaDto.fromJson(j)).toList();
\end{codigo}
    \end{column}
    \begin{column}{0.4\textwidth}
      \begin{conceito}[Dio]
        Cliente HTTP com URL base, tempos limite, \textbf{interceptadores}, cancelamento, \emph{upload} com progresso e decodificação automática de JSON.
      \end{conceito}
      \begin{conceito}[Interceptador]
        Código que roda em toda requisição, resposta ou erro: incluir token, registrar log, medir tempo, renovar sessão.
      \end{conceito}
    \end{column}
  \end{columns}
\end{frame}

\begin{frame}[fragile]{Erros e renovação de token}
  \begin{columns}[T]
    \begin{column}{0.52\textwidth}
\begin{codigo}[title=Traduzindo DioException]{Dart}
String mensagemDe(DioException e) => switch (e.type) {
  DioExceptionType.connectionTimeout ||
  DioExceptionType.receiveTimeout =>
      'O servidor demorou a responder.',
  DioExceptionType.connectionError =>
      'Sem conexão com a internet.',
  DioExceptionType.badResponse => switch (
      e.response?.statusCode) {
    401 => 'Sessão expirada.',
    404 => 'Recurso não encontrado.',
    _ => 'Erro no servidor.',
  },
  _ => 'Falha inesperada.',
};
\end{codigo}
    \end{column}
    \begin{column}{0.46\textwidth}
\begin{codigo}[title=Renovando o token após 401]{Dart}
QueuedInterceptorsWrapper(
  onError: (e, handler) async {
    if (e.response?.statusCode != 401) {
      return handler.next(e);
    }
    final ok = await auth.renovarToken();
    if (!ok) return handler.next(e);
    final repetida = await dio.fetch(
        e.requestOptions);
    handler.resolve(repetida);
  },
)
\end{codigo}
      \begin{dica}
        A versão \emph{Queued} processa um erro por vez: várias requisições com 401 disparam uma única renovação.
      \end{dica}
    \end{column}
  \end{columns}
\end{frame}

\begin{frame}[fragile]{Upload, download e cancelamento}
  \begin{columns}[T]
    \begin{column}{0.56\textwidth}
\begin{codigo}[title=Foto do equipamento com progresso]{Dart}
final cancelar = CancelToken();

Future<void> enviarFoto(int id, String caminho) =>
    dio.post('equipamentos/$id/foto',
      data: FormData.fromMap({
        'arquivo': await MultipartFile.fromFile(caminho,
            filename: 'foto.jpg'),
      }),
      cancelToken: cancelar,
      onSendProgress: (enviado, total) =>
          progresso.value = enviado / total,
    );

Future<void> baixarRelatorio(String destino) =>
    dio.download('relatorios/mensal.pdf', destino,
        cancelToken: cancelar);

// ao sair da tela
@override
void dispose() { cancelar.cancel(); super.dispose(); }
\end{codigo}
    \end{column}
    \begin{column}{0.42\textwidth}
      \begin{conceito}[\cod{multipart/form-data}]
        Formato usado para enviar arquivos junto com campos de formulário.
      \end{conceito}
      \begin{dica}
        Cancele requisições quando o usuário sai da tela: economiza dados e evita atualizar um estado que já não existe.
      \end{dica}
      \begin{atencao}
        \cod{MultipartFile.fromFile} usa \cod{dart:io}; na web, use \cod{MultipartFile.fromBytes}.
      \end{atencao}
    \end{column}
  \end{columns}
\end{frame}

% =============================================================================
\section{Robustez e tempo real}

\begin{frame}[fragile]{Falhas de rede são normais}
  \begin{columns}[T]
    \begin{column}{0.54\textwidth}
\begin{codigo}[title=utils/retentativa.dart]{Dart}
Future<T> comRetentativa<T>(Future<T> Function() acao,
    {int maximo = 3}) async {
  for (var i = 0; ; i++) {
    try {
      return await acao();
    } on DioException catch (e) {
      final transitoria =
          e.type != DioExceptionType.badResponse ||
          (e.response?.statusCode ?? 0) >= 500;
      if (!transitoria || i + 1 >= maximo) rethrow;
      // 0,5 s, 1 s, 2 s ...
      await Future.delayed(
          Duration(milliseconds: 500 * (1 << i)));
    }
  }
}
\end{codigo}
    \end{column}
    \begin{column}{0.44\textwidth}
      \begin{conceito}[Estratégia]
        \begin{itemize}
          \item Repita só erros \textbf{transitórios} (tempo limite, 5xx), nunca 400 ou 401.
          \item Espere cada vez mais (\emph{backoff} exponencial).
          \item O repositório converte a exceção em \cod{Result} (Módulo 08).
          \item Mostre dados do banco local enquanto a rede falha (Módulo 10).
        \end{itemize}
      \end{conceito}
      \begin{dica}
        \cod{connectivity\_plus} informa o tipo de rede (Wi-Fi, móvel), não se há internet de fato.
      \end{dica}
    \end{column}
  \end{columns}
\end{frame}

\begin{frame}[fragile]{Tempo real com WebSocket}
  \begin{columns}[T]
    \begin{column}{0.56\textwidth}
\begin{codigo}[title=Leituras ao vivo da bancada]{Dart}
import 'package:web_socket_channel/web_socket_channel.dart';

class LeiturasAoVivo {
  WebSocketChannel? _canal;

  Future<Stream<Leitura>> conectar(int bancada) async {
    final canal = WebSocketChannel.connect(
        Uri.parse('wss://api.labmonitor.dev/ws'));
    await canal.ready;                  // conexão aberta
    canal.sink.add(jsonEncode({'assinar': bancada}));
    _canal = canal;
    return canal.stream
        .map((m) => Leitura.fromJson(jsonDecode(m)));
  }

  Future<void> fechar() async => _canal?.sink.close();
}
\end{codigo}
    \end{column}
    \begin{column}{0.42\textwidth}
      \begin{conceito}[WebSocket]
        Conexão que fica aberta: o servidor envia dados assim que eles existem, sem o app perguntar de tempos em tempos (\emph{polling}).
      \end{conceito}
      \begin{dica}
        O resultado é um \cod{Stream}: use \cod{StreamBuilder} ou converta no view model (Módulo 06).
      \end{dica}
      \begin{atencao}
        Conexões caem. Reconecte com \emph{backoff} e feche o canal quando a tela sai.
      \end{atencao}
    \end{column}
  \end{columns}
\end{frame}

\begin{frame}[fragile]{Testando código de rede}
  \begin{columns}[T]
    \begin{column}{0.58\textwidth}
\begin{codigo}[title=test/api\_service\_test.dart]{Dart}
import 'package:http/http.dart' as http;
import 'package:http/testing.dart';

void main() {
  test('getBancadas converte o JSON', () async {
    final cliente = MockClient((req) async {
      expect(req.url.path, '/v1/bancadas');
      return http.Response(
          '[{"id":1,"nome":"B1","reservada_ate":null}]', 200);
    });
    final api = ApiService(client: cliente);
    final lista = await api.getBancadas();
    expect(lista.single.nome, 'B1');
  });

  test('status 500 lança ApiException', () async {
    final api = ApiService(client: MockClient(
        (_) async => http.Response('', 500)));
    expect(api.getBancadas(), throwsA(isA<ApiException>()));
  });
}
\end{codigo}
    \end{column}
    \begin{column}{0.4\textwidth}
      \begin{conceito}[Sem rede nos testes]
        \cod{MockClient} responde o que você quiser, de forma instantânea e determinística. Teste sucesso, erro, JSON malformado e tempo limite.
      \end{conceito}
      \begin{dica}
        Com Dio, substitua o \cod{httpClientAdapter} ou use o pacote \cod{http\_mock\_adapter}.
      \end{dica}
    \end{column}
  \end{columns}
\end{frame}

\begin{frame}{Segurança e escolha da biblioteca}
  \begin{columns}[T]
    \begin{column}{0.48\textwidth}
      \begin{atencao}[Segurança na rede]
        \begin{itemize}
          \item Só HTTPS; nada de desativar a validação de certificado.
          \item \textbf{Nunca} coloque segredos no app: o binário pode ser extraído. Chaves privadas ficam no backend.
          \item Tokens em armazenamento seguro (Módulo 10).
          \item Em apps sensíveis, considere \emph{certificate pinning}.
        \end{itemize}
      \end{atencao}
    \end{column}
    \begin{column}{0.5\textwidth}
      \scriptsize
      \begin{tabularx}{\linewidth}{@{}l L@{}}
        \toprule
        \cabecalho{Pacote} & \cabecalho{Indicado para} \\
        \midrule
        \cod{http} & simples, oficial do time Dart; poucas chamadas \\
        \textbf{Dio} & interceptadores, \emph{upload}, cancelamento; apps reais \\
        \cod{retrofit} & interface anotada gerando cliente Dio (como o Retrofit do Kotlin) \\
        \cod{web\_socket\_channel} & tempo real \\
        \bottomrule
      \end{tabularx}
      \begin{dica}
        O \textbf{DevTools} tem uma aba \emph{Network} que mostra cada requisição do app em execução.
      \end{dica}
    \end{column}
  \end{columns}
\end{frame}

% =============================================================================
\section{Fechamento}

\begin{frame}{Resumo}
  \begin{resumo}
    \begin{itemize}
      \item HTTP: método, URL, cabeçalhos e corpo; a resposta traz status e dados.
      \item \cod{http} é simples; Dio acrescenta interceptadores, cancelamento e progresso.
      \item Valide o JSON com \emph{patterns} ou gere o código com \cod{json\_serializable}.
      \item Trate erros por tipo, repita só falhas transitórias e devolva \cod{Result}.
      \item WebSocket para tempo real; \cod{MockClient} para testes sem rede.
    \end{itemize}
  \end{resumo}
\end{frame}

\begin{frame}{Glossário (1/2)}
  \begin{glossario}
    Protocolo & regras de formato e de ordem das mensagens entre computadores \\
    TCP/IP & pilha de protocolos da internet: aplicação, transporte, internet e enlace \\
    DNS & sistema que traduz nomes de domínio em endereços IP \\
    NAT & tradução de endereços que esconde o celular atrás de um IP público \\
    TCP & transporte com conexão, entrega garantida e ordem \\
    UDP & transporte sem conexão nem garantias, de baixa latência \\
    TLS & protocolo que cifra e autentica a conexão (o ``S'' do HTTPS) \\
    Latência (RTT) & tempo de ida e volta de uma mensagem \\
    HTTP/2 e HTTP/3 & vários pedidos na mesma conexão; o HTTP/3 roda sobre QUIC (UDP) \\
  \end{glossario}
\end{frame}

\begin{frame}{Glossário (2/2)}
  \begin{glossario}
    REST & estilo de API com recursos em URLs e métodos HTTP \\
    \emph{ETag} & identificador de versão que permite a resposta 304 sem corpo \\
    Serialização & conversão de objetos em bytes (JSON, Protocol Buffers) \\
    Idempotência & repetir a operação produz o mesmo resultado \\
    \cod{http} e Dio & pacotes de cliente HTTP do Dart; o Dio tem interceptadores e cancelamento \\
    \cod{json\_serializable} & gera o código de conversão entre JSON e classes Dart \\
    Interceptador & código que examina ou altera toda requisição e resposta \\
    WebSocket & canal aberto nos dois sentidos para tempo real \\
  \end{glossario}
\end{frame}

\begin{frame}{Questões de revisão}
  \begin{questoes}
    \item Descreva, camada por camada, o que acontece quando o app chama \cod{GET https://api.labmonitor.dev/v1/bancadas}.
    \item Por que o servidor não consegue abrir uma conexão até o celular? Que soluções existem?
    \item Por que, em apps, a latência costuma pesar mais que a largura de banda?
    \item Que três garantias o TLS oferece e como o app sabe que fala com o servidor certo?
    \item O que o HTTP/3 traz de especialmente útil para um celular em movimento?
    \item Explique como \emph{ETag} e a resposta 304 economizam dados.
    \item Que cuidados permitem mudar a API sem quebrar versões antigas do app?
    \item Uma requisição \cod{POST /reservas} deu tempo limite esgotado. É seguro repetir? Como tornar seguro?
  \end{questoes}
\end{frame}

\begin{frame}{Exercícios}
  \begin{exercicio}[Exercício 1 --- API pública]
    Consuma uma API pública (por exemplo, a de CEP ou de clima) com \cod{http} e mostre carregando, erro e dados.
  \end{exercicio}
  \begin{exercicio}[Exercício 2 --- Dio]
    Migre o serviço para Dio com interceptador de token e mensagens de erro traduzidas.
  \end{exercicio}
  \begin{exercicio}[Exercício 3 --- Testes]
    Escreva testes com \cod{MockClient} para 200, 404, 500 e JSON malformado.
  \end{exercicio}
  \begin{projeto}[LabMonitor --- Etapa 11]
    Sincronize bancadas e leituras com a API em Dio, salve no Drift (Módulo 10) e mostre as leituras ao vivo via WebSocket.
  \end{projeto}
\end{frame}

\begin{frame}{Para aprofundar}
  \small
  \begin{itemize}
    \item Rede no Flutter: \url{https://docs.flutter.dev/data-and-backend/networking}
    \item JSON e serialização: \url{https://docs.flutter.dev/data-and-backend/serialization/json}
    \item Pacote http: \url{https://pub.dev/packages/http} \quad Dio: \url{https://pub.dev/packages/dio}
    \item json\_serializable: \url{https://pub.dev/packages/json_serializable}
    \item WebSocket: \url{https://docs.flutter.dev/cookbook/networking/web-sockets}
  \end{itemize}
\end{frame}

\begin{frame}{Respostas das questões de revisão}
  \begin{respostas}
    \item O DNS traduz o nome em IP; o TCP abre a conexão na porta 443; o TLS combina as chaves e confere o certificado; o HTTP envia a requisição cifrada; tudo viaja em pacotes IP por Wi-Fi ou rede móvel.
    \item O celular fica atrás de NAT, com endereço privado, e não aceita conexões novas. Soluções: notificações \emph{push} (FCM, APNs) ou conexões abertas pelo próprio app (WebSocket, MQTT).
    \item As requisições são pequenas: o tempo de ida e volta (aperto de mão, DNS, TLS) domina o tempo de transferência.
    \item Confidencialidade, integridade e autenticidade. O certificado do servidor é assinado por uma autoridade em que o sistema confia e precisa valer para aquele domínio.
    \item Roda sobre QUIC: sem bloqueio na fila do TCP e com conexão que sobrevive à troca de rede, como do Wi-Fi para o 4G.
    \item O app envia a versão que tem (\cod{If-None-Match}); se nada mudou, o servidor responde 304 sem corpo e o app reutiliza a cópia local.
    \item Só acrescentar campos, ignorar campos desconhecidos (o padrão do \cod{json\_serializable}), ter valores padrão e criar \cod{/v2} para mudanças incompatíveis.
    \item Não, pois \cod{POST} não é idempotente e pode criar duas reservas. Envie uma chave de idempotência, para o servidor reconhecer a repetição.
  \end{respostas}
\end{frame}

\encerramento{Módulo 12 --- Firebase com FlutterFire: autenticação, Firestore, Storage, mensagens e Crashlytics.}

\end{document}
